% doc/cc_completeness.tex -- completeness of the connected-cluster (CC) search orders of dist_m4ri
% Build: pdflatex cc_completeness.tex (twice, for the references)
\documentclass[11pt]{article}
\usepackage[margin=1in]{geometry}
\usepackage{amsmath,amssymb,amsthm,mathtools}
\usepackage{enumitem}
\usepackage[colorlinks=true,linkcolor=blue,citecolor=blue,urlcolor=blue]{hyperref}

\newtheorem{theorem}{Theorem}
\newtheorem{lemma}[theorem]{Lemma}
\newtheorem{proposition}[theorem]{Proposition}
\newtheorem{corollary}[theorem]{Corollary}
\theoremstyle{definition}
\newtheorem{definition}[theorem]{Definition}
\newtheorem{remark}[theorem]{Remark}

\newcommand{\Ftwo}{\mathbb{F}_2}
\DeclareMathOperator{\supp}{supp}
\DeclareMathOperator{\rank}{rank}
\newcommand{\Tr}{\mathcal{T}}   % trivial codewords
\newcommand{\Kc}{\mathcal{K}}   % completions
\newcommand{\Wr}{\mathcal{W}}   % reported weights
\newcommand{\ttt}[1]{\texttt{#1}}

\begin{document}

\title{Completeness of the connected-cluster search orders in \ttt{dist\_m4ri}:\\
any unsatisfied row, forbidden sets, and admissible bounds}
\author{Notes for \ttt{dist-m4ri}}
\date{October 2026}
\maketitle

\begin{abstract}
The connected-cluster (CC) algorithm of Dumer, Kovalev, and Pryadko~\cite{DKP17} certifies lower bounds on the
distance by enumerating irreducible clusters: starting from a column of the parity-check matrix, a cluster is extended,
one column at a time, by a column of an unsatisfied row. Up to version 0.14.7, \ttt{dist\_m4ri} always extends a
cluster by the columns of the \emph{first} unsatisfied row in row order. We formulate the modified order of the
branch-and-bound search of Ref.~\cite{LOMKZ26} (App.~B), used in Ref.~\cite{Lee26} (App.~C.4): branching on an
\emph{arbitrary} unsatisfied row, in particular on the one with the fewest available columns; a \emph{forbidden set}
which excludes the columns of the earlier sibling branches; and pruning by \emph{admissible lower bounds} on the number
of missing columns. Theorem~\ref{thm:main} states that the search reports every irreducible codeword of weight at most
$W$ which contains the start column and avoids the initial forbidden set, and that it reports it exactly once if the
forbidden set is used. The current order is the special case without the forbidden set, where every irreducible
codeword is reported at least once. The corollaries cover the certification of the distance by increasing weights, the
hit counts, start lists, orbit exclusion under a group of code automorphisms, the hash table used for the last two
extensions of a cluster, and the work units of the multithreaded search.
\end{abstract}

\section{Setting and notation}\label{sec:setting}

Let $H\in\Ftwo^{m\times n}$ be a parity-check matrix with columns $h_1,\dots,h_n\in\Ftwo^m$. The columns are the
variable nodes of a code (or the error mechanisms of a detector error model, DEM), and the rows are its checks (or
detectors). We write $[n]=\{1,\dots,n\}$ with its natural order. For a row $r\in[m]$ let
$N(r)=\{j\in[n]:H_{rj}=1\}$, and for a column $j$ let $R(j)=\{r\in[m]:H_{rj}=1\}$; $\Delta=\max_j|R(j)|$ is the
maximum column weight. A set $S\subseteq[n]$ is identified with its indicator vector $1_S\in\Ftwo^n$. The
\emph{syndrome} of $S$ is
\[
  \sigma(S)=\sum_{j\in S}h_j=H\,1_S\in\Ftwo^m ,
\]
and $U(S)=\supp\sigma(S)=\{r\in[m]:|N(r)\cap S|\text{ is odd}\}$ is the set of \emph{unsatisfied rows} of $S$. For
disjoint sets $A,B\subseteq[n]$,
\begin{equation}\label{eq:additive}
  \sigma(A\cup B)=\sigma(A)+\sigma(B).
\end{equation}

A set $C\subseteq[n]$ is a \emph{codeword} if $\sigma(C)=0$, i.e., $1_C\in\ker H$. Let $\Tr\subseteq\ker H$ be a
linear subspace, the \emph{trivial} codewords. A codeword $C$ is \emph{nontrivial} if $1_C\notin\Tr$, and the
\emph{distance} is $d=\min\{|C|: C\text{ a nontrivial codeword}\}$. For a classical code $\Tr=\{0\}$; for one sector
of a CSS code with $H=H_X$, $\Tr$ is the row space of $H_Z$; for a DEM with the observable matrix $L$,
$\Tr=\ker H\cap\ker L$. In \ttt{dist\_m4ri} nontriviality is tested as $L\,1_C\neq0$ with a matrix $L$ such that
$\ker H\cap\ker L=\Tr$. Only the linearity of $\Tr$ is used below.

\section{Irreducible codewords}\label{sec:irreducible}

\begin{definition}[irreducible codewords~\cite{DKP17}]\label{def:irreducible}
A nonempty codeword $C$ is \emph{reducible} if $C=A\cup B$ for two disjoint nonempty codewords $A$ and $B$; otherwise
it is \emph{irreducible}.
\end{definition}

\begin{lemma}[equivalent forms]\label{lem:forms}
Let $C$ be a nonempty codeword, and let $H_C$ be the $m\times|C|$ submatrix of the columns of $H$ in $C$.
\begin{enumerate}[label=(\alph*)]
  \item $C$ is reducible if and only if some nonempty proper subset $S\subsetneq C$ is a codeword.
  \item $\rank H_C\le|C|-1$, and $C$ is reducible if and only if $\rank H_C\le|C|-2$. Thus $C$ is irreducible if and
    only if $\rank H_C=|C|-1$.
  \item If $C$ is nontrivial, then $C$ is reducible if and only if $C=A\cup B$ with a nontrivial codeword $A$ and a
    nonempty codeword $B$ disjoint from $A$ (``a codeword and a zero-syndrome vector with disjoint supports'').
\end{enumerate}
\end{lemma}

\begin{proof}
(a) If $S\subsetneq C$ is a nonempty codeword, then $B=C\setminus S$ is nonempty and $\sigma(B)=\sigma(C)+\sigma(S)=0$
by~\eqref{eq:additive}, so $C=S\cup B$ is reducible. The converse is immediate from the definition.

(b) The all-ones vector $\mathbf{1}\in\Ftwo^{|C|}$ lies in $\ker H_C$ since $\sigma(C)=0$, so
$\dim\ker H_C=|C|-\rank H_C\ge1$. If $\dim\ker H_C\ge2$, there is $x\in\ker H_C\setminus\{0,\mathbf{1}\}$; its support
is a nonempty proper subset of $C$ with zero syndrome. Conversely, the indicator vector of a nonempty proper subset of
$C$ with zero syndrome lies in $\ker H_C\setminus\{0,\mathbf{1}\}$, so that $\dim\ker H_C\ge2$. Combine with~(a).

(c) If $C=A\cup B$ with disjoint codewords, then $1_C=1_A+1_B$; since $\Tr$ is linear and $1_C\notin\Tr$, at least one
of $1_A$ and $1_B$ is not in $\Tr$; call that set $A$. The converse holds since a nontrivial codeword is nonempty.
\end{proof}

\begin{lemma}[minimum-weight codewords are irreducible]\label{lem:minweight}
A reducible nontrivial codeword $C$ contains a nontrivial codeword of weight at most $|C|-1$. In particular, every
nontrivial codeword of weight $d$ is irreducible.
\end{lemma}

\begin{proof}
By Lemma~\ref{lem:forms}(c), $C=A\cup B$ with $A$ nontrivial and $B$ nonempty and disjoint from $A$, so
$|A|\le|C|-1$. A nontrivial codeword of weight $d$ cannot contain a nontrivial codeword of smaller weight.
\end{proof}

\begin{lemma}[parity]\label{lem:parity}
Let $C$ be a codeword, $S\subseteq C$, and $r\in U(S)$. Then $|N(r)\cap(C\setminus S)|$ is odd; in particular,
$N(r)\cap(C\setminus S)\neq\emptyset$.
\end{lemma}

\begin{proof}
By~\eqref{eq:additive}, $\sigma(C\setminus S)=\sigma(C)+\sigma(S)=\sigma(S)$, whose entry in the row $r$ is $1$.
\end{proof}

\begin{remark}[connectivity]\label{rem:connected}
Call two columns adjacent if they share a row. An irreducible codeword is connected: the columns of a connected
component of a codeword share no row with the other columns, so each component is itself a codeword. The clusters of
the search below are connected as well, since each added column lies in a row $r\in U(S)$, and $|N(r)\cap S|$ is odd,
hence positive. The converse fails: the disjoint union of two codewords which share a row (each meeting it an even
number of times) is connected and reducible. Such codewords are found, e.g., in the rounds $w\ge d+3$ of the CC search
on circuit DEMs, which have many trivial codewords of weight~3. Therefore a test of connectivity does not detect them,
while the rank test of Lemma~\ref{lem:forms}(b) does.
\end{remark}

\section{The search}\label{sec:search}

A \emph{node} is a pair $(S,F)$ of disjoint subsets of $[n]$ with $S\neq\emptyset$ and $\sigma(S)\neq0$; $S$ is the
\emph{cluster} and $F$ the \emph{forbidden set}. A column $u$ is \emph{available} at $(S,F)$ if $u\notin S\cup F$, and
\[
  A_r(S,F)=N(r)\setminus(S\cup F)
\]
is the set of available columns of the row $r$.

\begin{definition}[completions, admissible bounds]\label{def:admissible}
The \emph{completions} of a node $(S,F)$ are the sets
\[
  \Kc(S,F)=\{T\subseteq[n]\setminus(S\cup F):\ \sigma(T)=\sigma(S)\},
\]
i.e., the sets $T$ of available columns for which $S\cup T$ is a codeword, and
$\lambda(S,F)=\min\{|T|:T\in\Kc(S,F)\}$, with $\lambda(S,F)=\infty$ if $\Kc(S,F)=\emptyset$. A function $\beta$ from
the nodes to $\mathbb{Z}_{\ge0}\cup\{\infty\}$ is an \emph{admissible bound} if $\beta(S,F)\le\lambda(S,F)$ for every
node $(S,F)$.
\end{definition}

The pointwise maximum of admissible bounds is admissible, and so is a bound which is evaluated only at some nodes and
set to $0$ at the others (e.g., only where it can prune).

\begin{lemma}[admissible bounds]\label{lem:bounds}
For a node $(S,F)$ let $U=U(S)$, and for an available column $u$ let $c_u=|R(u)\cap U|$ be the number of unsatisfied
rows which contain $u$. The following functions are admissible bounds:
\begin{enumerate}[label=(\alph*)]
  \item \emph{dead end}: $\beta_0=\infty$ if $A_r(S,F)=\emptyset$ for some $r\in U$, and $\beta_0=0$ otherwise;
  \item \emph{syndrome weight}: $\beta_1=\lceil|U|/\Delta\rceil$;
  \item \emph{packing}: $\beta_2=|P|$ for any set $P\subseteq U$ of unsatisfied rows whose sets $A_r(S,F)$, $r\in P$,
    are pairwise disjoint (e.g., a packing found greedily, the rows with the fewest available columns first);
  \item \emph{coverage}: $\beta_3=$ the least $t$ such that the $t$ largest of the values $c_u$ of the available
    columns $u$ sum to at least $|U|$ ($\beta_3=\infty$ if all of them sum to less than $|U|$);
  \item \emph{fractional cover}: $\beta_4=\bigl\lceil\sum_{r\in U}1/\mu_r\bigr\rceil$ with
    $\mu_r=\max_{u\in A_r(S,F)}c_u$ ($\beta_4=\infty$ if $A_r(S,F)=\emptyset$ for some $r\in U$).
\end{enumerate}
\end{lemma}

\begin{proof}
If $\Kc(S,F)=\emptyset$, then $\lambda(S,F)=\infty$ and there is nothing to prove. Let $T\in\Kc(S,F)$. For every
$r\in U$, the entry $r$ of $\sigma(T)=\sigma(S)$ is $1$, so $|N(r)\cap T|$ is odd and at least $1$, and
$N(r)\cap T\subseteq A_r(S,F)$ since $T$ contains no column of $S\cup F$.
(a) Hence $A_r(S,F)\neq\emptyset$ for every $r\in U$.
(c) Choosing a column $u_r\in N(r)\cap T\subseteq A_r(S,F)$ for each $r\in P$ gives $|P|$ distinct columns of $T$.
(b), (d) Counting the incidences between $T$ and $U$ in two ways,
\[
  |U|\le\sum_{r\in U}|N(r)\cap T|=\sum_{u\in T}c_u\le(\text{the sum of the }|T|\text{ largest values }c_u)\le\Delta|T|.
\]
(e) For an available column $u$ and a row $r\in U\cap R(u)$, $u\in A_r(S,F)$, so that $\mu_r\ge c_u$; hence
$\sum_{r\in U\cap R(u)}1/\mu_r\le c_u/c_u=1$ (and the sum is empty if $c_u=0$). Therefore
\[
  |T|\ \ge\ \sum_{u\in T}\ \sum_{r\in U\cap R(u)}\frac1{\mu_r}\ =\ \sum_{r\in U}\frac{|N(r)\cap T|}{\mu_r}\
  \ge\ \sum_{r\in U}\frac1{\mu_r},
\]
and $|T|$ is an integer.
\end{proof}

\paragraph{The algorithm.} The search has the following parameters:
\begin{itemize}
  \item a weight limit $W\ge1$ and a set $\Wr\subseteq\{1,\dots,W\}$ of \emph{reported weights} ($\Wr=\{W\}$ in the
    round $W$ of \ttt{dist\_m4ri});
  \item a \emph{choice rule} $\rho$, which assigns to each node $(S,F)$ an unsatisfied row $\rho(S,F)\in U(S)$;
  \item for each node, an ordering $(v_1,\dots,v_q)$ of the set $A_{\rho(S,F)}(S,F)$ (e.g., by column index);
  \item a flag $\varepsilon\in\{0,1\}$: \emph{sibling exclusion} (the forbidden set) on or off;
  \item an admissible bound $\beta$.
\end{itemize}
The choice rule and the ordering must be functions of the node $(S,F)$; they may depend on $S$ and $F$ in any way,
e.g., through the numbers $|A_r(S,F)|$.

\medskip
\noindent\textup{\textsc{Search}}$(S,F)$ \quad (called on nodes $(S,F)$ with $|S|<W$):
\begin{enumerate}[label=\arabic*.,leftmargin=2.5em]
  \item (Pruning.) If $|S|+\beta(S,F)>W$, return.
  \item (Branching row.) Let $r=\rho(S,F)$, and let $(v_1,\dots,v_q)$ be the ordering of $A_r(S,F)$.
  \item For $i=1,\dots,q$:
  \begin{enumerate}[label=(\alph*)]
    \item $S_i=S\cup\{v_i\}$; $F_i=F\cup\{v_1,\dots,v_{i-1}\}$ if $\varepsilon=1$, and $F_i=F$ if $\varepsilon=0$;
    \item if $\sigma(S_i)=0$: report $S_i$ if $|S_i|\in\Wr$, and do not search below $S_i$;
    \item otherwise, if $|S_i|<W$: call \textup{\textsc{Search}}$(S_i,F_i)$.
  \end{enumerate}
\end{enumerate}
The \emph{search from} a start column $v$ with an initial forbidden set $F_0\not\ni v$ reports $\{v\}$ if
$\sigma(\{v\})=0$ and $1\in\Wr$; if $\sigma(\{v\})\neq0$ and $W\ge2$, it calls \textup{\textsc{Search}}$(\{v\},F_0)$.

In step~3(a), the columns $v_1,\dots,v_{i-1}$ are forbidden in the branch $i$ whatever happened in their own branches
(searched, pruned, reported, or ended at a zero syndrome).

\paragraph{Instances.}
\begin{itemize}
  \item \emph{The current order} (\ttt{cc\_walk()} and \ttt{start\_CC\_recurs\_mt()} in \ttt{dist\_m4ri}~0.14.7,
    the irreducible-cluster principle of~\cite{DKP17}): $\rho_{\rm first}(S,F)=\min U(S)$, the first unsatisfied row
    in the row order of $H$ (rows sorted by weight with \ttt{sort\_rows=1}); $\varepsilon=0$; $\beta=\beta_1$ (tested
    before a child is searched, which is the same as step~1 of the child); the columns in increasing order; and
    $F_0=\{u\in[n]:u<v\}$, i.e., clusters grow by columns larger than the start column.
  \item \emph{The branch and bound} of~\cite{LOMKZ26} (\ttt{cc\_prune=1}): $\rho_{\min}(S,F)$ is an unsatisfied row
    with the fewest available columns, $|A_r(S,F)|$ minimal (ties broken by the row order); $\varepsilon=1$; and
    $\beta=\max(\beta_0,\beta_1,\beta_2)$, with $\beta_2$ computed greedily and only where it can prune
    ($|U(S)|>W-|S|$) and $W-|S|\ge3$. With \ttt{cc\_prune=2}, also $\beta_3$ and $\beta_4$.
  \item \emph{Unlimited clusters} from a start list (\ttt{start=}): $F_0=\emptyset$, or the earlier start columns
    (Proposition~\ref{prop:start}), or the earlier orbits (Proposition~\ref{prop:orbits}).
\end{itemize}

\section{Main theorem}\label{sec:main}

\begin{theorem}[completeness and uniqueness]\label{thm:main}
Fix the parameters of the search, a start column $v$, and an initial forbidden set $F_0\not\ni v$. Let $C$ be an
irreducible codeword with $|C|\in\Wr$ (so that $|C|\le W$), $v\in C$, and $C\cap F_0=\emptyset$. Then:
\begin{enumerate}[label=(\roman*)]
  \item the search from $v$ reports $C$;
  \item if $\varepsilon=1$, the search from $v$ reports no set more than once; in particular, it reports $C$ exactly
    once;
  \item every set $R$ reported by the search from $v$ is a codeword with $|R|\in\Wr$, $v\in R$, and
    $R\cap F_0=\emptyset$.
\end{enumerate}
\end{theorem}

\begin{proof}
\emph{(iii) Soundness.} Along the calls, the cluster grows by one available column at a time, the forbidden set does
not decrease, and $S_i\cap F_i=\emptyset$: indeed $v_i\notin S\cup F$, and $v_1,\dots,v_{i-1}$ are available at
$(S,F)$, hence not in $S$, and distinct from $v_i$. A reported set is $\{v\}$ or a set $S_i$ with $\sigma(S_i)=0$ and
$|S_i|\in\Wr$. It contains $v$, and none of its columns is in $F_0$: each was available when it was added, i.e., not
in the current forbidden set, which contains $F_0$, and $v\notin F_0$.

\emph{(i) Existence.} If $|C|=1$, then $C=\{v\}$ and $\sigma(\{v\})=0$, and $\{v\}$ is reported. Let $|C|\ge2$. We
show by induction on $t=0,1,\dots$ that either $C$ has been reported, or \textup{\textsc{Search}} is called on a node
$(S_t,F_t)$ with
\begin{equation}\label{eq:invariant}
  v\in S_t\subsetneq C,\qquad |S_t|=t+1,\qquad F_t\cap C=\emptyset .
\end{equation}
For $t=0$: $\{v\}\subsetneq C$ is nonempty, so $\sigma(\{v\})\neq0$ by Lemma~\ref{lem:forms}(a), and $W\ge|C|\ge2$;
hence \textup{\textsc{Search}}$(\{v\},F_0)$ is called, and $(S_0,F_0)=(\{v\},F_0)$ satisfies~\eqref{eq:invariant}.

For the induction step, let $(S,F)=(S_t,F_t)$ satisfy~\eqref{eq:invariant}, and let $T=C\setminus S\neq\emptyset$.
\begin{itemize}
  \item $T$ is a completion of $(S,F)$: $T\cap S=\emptyset$, $T\cap F\subseteq C\cap F=\emptyset$, and
    $\sigma(T)=\sigma(C)+\sigma(S)=\sigma(S)$ by~\eqref{eq:additive}. Hence
    $\beta(S,F)\le\lambda(S,F)\le|T|=|C|-|S|\le W-|S|$, and step~1 does not return.
  \item The row $r=\rho(S,F)$ is unsatisfied, so $N(r)\cap T\neq\emptyset$ by Lemma~\ref{lem:parity}. The columns of
    $T$ are available, so $N(r)\cap T\subseteq A_r(S,F)=\{v_1,\dots,v_q\}$. Let $j$ be the least index with $v_j\in C$.
    Then $S_j=S\cup\{v_j\}\subseteq C$, and $F_j\cap C=\emptyset$, since $F_j\subseteq F\cup\{v_1,\dots,v_{j-1}\}$ and
    $v_1,\dots,v_{j-1}\notin C$ by the choice of $j$.
  \item If $S_j=C$, then $\sigma(S_j)=0$ and $|S_j|=|C|\in\Wr$, so $C$ is reported in step~3(b).
  \item If $S_j\subsetneq C$, then $\sigma(S_j)\neq0$ by Lemma~\ref{lem:forms}(a), since $C$ is irreducible, and
    $|S_j|<|C|\le W$; hence \textup{\textsc{Search}}$(S_j,F_j)$ is called in step~3(c), and
    $(S_{t+1},F_{t+1})=(S_j,F_j)$ satisfies~\eqref{eq:invariant}.
\end{itemize}
Since $|S_t|=t+1<|C|$ in~\eqref{eq:invariant}, the induction ends after at most $|C|-1$ steps, with the report of $C$.

\emph{(ii) Uniqueness.} Let $\varepsilon=1$. A report is produced by a \emph{report event}: a call
\textup{\textsc{Search}}$(S,F)$ together with an index $i$ in step~3(b), or the start step, which reports $\{v\}$
and calls nothing else. A report event determines the sequence of branch indices taken from the start. Let
$e\neq e'$ be two report events of the search from $v$, and let \textup{\textsc{Search}}$(S,F)$ be the last call
common to both sequences, where they continue with the indices $i$ and $i'$. Then $i\neq i'$: if both events occur in
this call, they differ in the index; and if $e$ reports $S_i$ in this call, then $e'$ cannot lie below the branch $i$,
since nothing is searched below a set with zero syndrome. Suppose $i<i'$. The set reported by $e$ contains
$S_i\ni v_i$, since a cluster only grows along the calls. In the branch $i'$ and below it, $v_i\in F_{i'}$ is
forbidden, since the forbidden sets do not decrease, so $v_i$ is never added there, and $v_i\notin S\cup\{v_{i'}\}$.
Hence $v_i$ is not in the set reported by $e'$, and the two reported sets differ.
\end{proof}

\begin{remark}\label{rem:mult}
Without the forbidden set ($\varepsilon=0$), a codeword can be reported several times: if it meets the branching row
of a node in three columns which are not yet in the cluster, each of the three branches can lead to it. For instance,
the current order records 2084 hits for the 1884 codewords of weight 12 of the gross code $[[144,12,12]]$, and the
branch and bound records each of them once.
\end{remark}

\begin{remark}[what the proof uses]\label{rem:uses}
The proof uses only Lemma~\ref{lem:parity} for the chosen row, the admissibility of $\beta$, the irreducibility of
$C$, and the invariant $F_t\cap C=\emptyset$. Hence Theorem~\ref{thm:main} holds for every choice rule (the first
unsatisfied row, the row with the fewest available columns, or any other), for every ordering of the available
columns, and for every admissible bound. These choices change only the shape and the size of the search tree.
\end{remark}

\begin{remark}[branching factor]\label{rem:branching}
At every node, the number of branches is $q=|A_{\rho(S,F)}(S,F)|\le|N(\rho(S,F))|-1\le\max_r|N(r)|-1$, since the
chosen row is unsatisfied and therefore contains a column of $S$; this bound, and every bound on the size of the
search tree derived from it, holds for every choice rule. The rule $\rho_{\min}$ branches at each node on
$\min_{r\in U(S)}|A_r(S,F)|\le|A_{\rho_{\rm first}(S,F)}(S,F)|$ columns, i.e., never on more columns than the first
unsatisfied row offers at the same node, and it detects a dead end (the minimum $0$) at once. The total number of calls
is not monotone in these local choices, since another row leads to other clusters below the node. In a prototype
(October 2026, see \ttt{todo.md}), $\rho_{\min}$ alone (with $\varepsilon=0$ and $\beta=\beta_1$) reduced the number of
calls by factors of 1.05--1.7 on seven codes and DEMs, and increased it on none; together with the forbidden set and
the packing bound $\beta_2$, by factors of 1.7--66.
\end{remark}

\section{Consequences}\label{sec:consequences}

\begin{corollary}[all start columns]\label{cor:roots}
Run the search from every start column $v\in[n]$ with $F_0(v)=\{u\in[n]:u<v\}$. Every irreducible codeword $C$ with
$|C|\in\Wr$ is reported from $v=\min C$, at least once, and exactly once if $\varepsilon=1$; it is not reported from
any other start column.
\end{corollary}

\begin{proof}
From $v=\min C$: $C\cap F_0(v)=\emptyset$, and Theorem~\ref{thm:main} applies. From $v\notin C$: every set reported
from $v$ contains $v$. From $v\in C$ with $v>\min C$: $\min C\in F_0(v)$, while the sets reported from $v$ avoid
$F_0(v)$.
\end{proof}

With $\Wr=\{W\}$, only the start columns $v\le n-W+1$ need to be searched, since a set of weight $W$ with the
smallest column $v$ has $W-1$ columns larger than $v$. The expert options \ttt{cbeg} and \ttt{cend} restrict the start
columns to an interval; then exactly the irreducible codewords $C$ with $\min C$ in the interval are reported.

\begin{corollary}[distance by increasing weights]\label{cor:ladder}
Let $w_0\ge1$ be such that there is no nontrivial codeword of weight less than $w_0$. For $w=w_0,w_0+1,\dots$, run the
searches of Corollary~\ref{cor:roots} with $W=w$ and $\Wr=\{w\}$ (the \emph{round} $w$), and test every reported set
for nontriviality.
\begin{enumerate}[label=(\alph*)]
  \item If the rounds $w_0,\dots,W$ report no nontrivial codeword, then $d>W$.
  \item If $w^*$ is the first round which reports a nontrivial codeword, then $d=w^*$, and every nontrivial codeword
    of weight $d$ is reported in the round $w^*$, exactly once if $\varepsilon=1$.
\end{enumerate}
\end{corollary}

\begin{proof}
By Lemma~\ref{lem:minweight}, every nontrivial codeword of weight $d$ is irreducible, and $d\ge w_0$ by assumption;
by Corollary~\ref{cor:roots}, it is reported in the round $d$. (a) Hence $d\notin\{w_0,\dots,W\}$, i.e., $d>W$.
(b) The nontrivial codeword reported in the round $w^*$ gives $d\le w^*$, and $d<w^*$ is excluded since the round $d$
would report a nontrivial codeword. The rest follows from Corollary~\ref{cor:roots}.
\end{proof}

In \ttt{dist\_m4ri}, $w_0$ is the supplied \ttt{dmin} (default $1$). The option \ttt{noscan=1} runs the single round
$w=\ttt{wmax}$, which is certified only if $w_0=\ttt{wmax}$ is supplied: a round $w$ reports only codewords of weight
$w$, and it does not search below a zero-syndrome cluster of smaller weight (step~3(b)), which can lose only reducible
codewords of weight $w$ (Lemma~\ref{lem:forms}(a)), and no irreducible ones.

\begin{remark}[reducible codewords and the rank filter]\label{rem:reducible}
In a round $w>d$, a reducible codeword of weight $w$ is reported only if the search reaches it along a branch on which
no cluster before the last one has zero syndrome; whether this happens depends on the choice rule, the orderings, and
$\varepsilon$. If every reported set $R$ is accepted only if $\rank H_R=|R|-1$ (Lemma~\ref{lem:forms}(b)), then the
accepted sets of the round $w$ are exactly the irreducible codewords of weight $w$ (Corollary~\ref{cor:roots}),
independently of the choice rule and the orderings, and with $\varepsilon=1$ each of them is accepted exactly once:
the CC search contributes exactly one hit to each codeword it finds. The test costs a Gaussian elimination of $|R|$
columns restricted to the at most $|R|\Delta$ rows they touch, and it is applied only to the reported sets.
\end{remark}

\begin{proposition}[start lists]\label{prop:start}
Let $Q=\{q_1<\dots<q_k\}\subseteq[n]$ be a list of start columns, searched with unlimited clusters, and let $C$ be an
irreducible codeword with $|C|\in\Wr$ and $C\cap Q\neq\emptyset$.
\begin{enumerate}[label=(\alph*)]
  \item With $F_0(q_j)=\emptyset$ for all $j$, $C$ is reported from every $q_j\in C$, and from no other start column;
    with $\varepsilon=1$, it is reported exactly $|C\cap Q|$ times.
  \item With $F_0(q_j)=\{q_1,\dots,q_{j-1}\}$, $C$ is reported from $q_j$ with $j=\min\{i:q_i\in C\}$, and from no
    other start column; with $\varepsilon=1$, it is reported exactly once.
\end{enumerate}
Thus both variants report the same irreducible codewords, those which meet $Q$; forbidding the earlier start columns
changes only the multiplicities. In particular, every conclusion about $d$ drawn from~(a), e.g., under a symmetry
which maps every nontrivial codeword of weight $d$ to one which meets $Q$, holds for~(b).
\end{proposition}

\begin{proof}
(a) Theorem~\ref{thm:main} with $F_0=\emptyset$ for each $q_j\in C$; the sets reported from $q_j$ contain $q_j$.
(b) For $j=\min\{i:q_i\in C\}$, $C\cap F_0(q_j)=\emptyset$, and Theorem~\ref{thm:main} applies. The sets reported from
$q_i\notin C$ contain $q_i$; and if $q_i\in C$ with $i>j$, then $q_j\in C\cap F_0(q_i)$, while the sets reported from
$q_i$ avoid $F_0(q_i)$.
\end{proof}

\begin{proposition}[orbit exclusion]\label{prop:orbits}
Let $\Gamma$ be a group of permutations of $[n]$ which act on $\Ftwo^n$ by permuting the coordinates, with
$g(\ker H)=\ker H$ and $g(\Tr)=\Tr$ for all $g\in\Gamma$. Let $O_1,\dots,O_k$ be the orbits of $\Gamma$ on $[n]$, let
$q_j\in O_j$, and search from each $q_j$ with unlimited clusters and $F_0(q_j)=O_1\cup\dots\cup O_{j-1}$.
\begin{enumerate}[label=(\alph*)]
  \item For every irreducible codeword $C$ with $|C|\in\Wr$, there is $g\in\Gamma$ such that $g(C)$ is reported;
    $g(C)$ is an irreducible codeword of weight $|C|$, and it is nontrivial if and only if $C$ is.
  \item Every reported set is reported from exactly one start column, and with $\varepsilon=1$ exactly once.
\end{enumerate}
Consequently, Corollary~\ref{cor:ladder}(a) holds for these searches, and so does~(b) with ``every nontrivial codeword
of weight $d$ is reported'' replaced by ``an image under $\Gamma$ of every nontrivial codeword of weight $d$ is
reported''.
\end{proposition}

\begin{proof}
For $A\subseteq[n]$ and $g\in\Gamma$, $1_{g(A)}=g(1_A)$, so $g(A)$ is a codeword if and only if $A$ is. Hence $g$
maps the nonempty proper subsets of $C$ with zero syndrome to those of $g(C)$, and $g(C)$ is irreducible by
Lemma~\ref{lem:forms}(a); since $g(\Tr)=\Tr$, it is nontrivial if and only if $C$ is.
(a) Let $j=\min\{i:C\cap O_i\neq\emptyset\}$, $u\in C\cap O_j$, and $g\in\Gamma$ with $g(u)=q_j$. Since
$g(O_i)=O_i$, we have $g(C)\cap O_i=g(C\cap O_i)=\emptyset$ for $i<j$, so $g(C)\cap F_0(q_j)=\emptyset$ and
$q_j\in g(C)$; Theorem~\ref{thm:main} applies.
(b) A set reported from $q_i$ contains $q_i$ and avoids $O_1,\dots,O_{i-1}$, so that $i=\min\{i':R\cap O_{i'}\neq
\emptyset\}$ is determined by the reported set $R$; within one start column, use Theorem~\ref{thm:main}(ii).
For the consequence, repeat the proof of Corollary~\ref{cor:ladder} with an image $g(C)$ of a nontrivial codeword $C$
of weight $d$, which is again nontrivial, irreducible, and of weight $d$.
\end{proof}

For a CSS code, the conditions on $\Gamma$ hold if every $g\in\Gamma$ maps the row spaces of $H_X$ and $H_Z$ to
themselves; they can be verified for each generator by the rank tests $\rank[H;Hg]=\rank H$, and the same for $H_Z$
(or for $[H;L]$ when only $L$ is given), where $Hg$ is $H$ with its columns permuted by $g$. With trivial orbits (the
trivial group, $k=n$, $q_j=j$), Proposition~\ref{prop:orbits} is Corollary~\ref{cor:roots}.

\section{Implementation details covered by the theorem}\label{sec:impl}

\begin{proposition}[the last two extensions by a hash table]\label{prop:hash}
Let $\Wr=\{W\}$, and let $(S,F)$ be a node with $|S|=W-2$ which is not pruned in step~1; let $r=\rho(S,F)$ and let
$(v_1,\dots,v_q)$ be the ordering of $A_r(S,F)$. The sets reported by \textup{\textsc{Search}}$(S,F)$, including its
descendants, are exactly the sets $S\cup\{v_i,u\}$ for the pairs $(i,u)$ with $i\in[q]$ and
\begin{equation}\label{eq:pair}
  u\notin S\cup F\cup\{v_i\},\qquad h_u=\sigma(S)+h_{v_i}\neq0 ,
\end{equation}
each exactly once, and different pairs give different sets, for $\varepsilon=0$ and for $\varepsilon=1$. Likewise,
for a node $(S,F)$ with $|S|=W-1$ which is not pruned, the reported sets are the sets $S\cup\{u\}$ with
$u\notin S\cup F$ and $h_u=\sigma(S)$.
\end{proposition}

\begin{proof}
Consider the branch $i$: $S_i=S\cup\{v_i\}$ and $F_i\subseteq F\cup\{v_1,\dots,v_{i-1}\}$. If $\sigma(S_i)=0$, the
branch reports nothing ($|S_i|=W-1\notin\Wr$) and searches nothing, and no $u$ satisfies~\eqref{eq:pair}. Otherwise
\textup{\textsc{Search}}$(S_i,F_i)$ is called. If it is pruned, then $\lambda(S_i,F_i)\ge\beta(S_i,F_i)>1$, so no
available $u$ has $h_u=\sigma(S_i)$, since $\{u\}$ would be a completion of size $1$; nothing is lost. If it is not
pruned, it branches over $A_{r'}(S_i,F_i)$ with $r'=\rho(S_i,F_i)\in U(S_i)$, and reports $S_i\cup\{u\}$ exactly for
the $u$ in this set with $\sigma(S_i\cup\{u\})=0$, i.e., $h_u=\sigma(S_i)$. Since $r'\in\supp\sigma(S_i)$, every $u$
with
$h_u=\sigma(S_i)$ lies in $N(r')$; so these are all the columns $u\notin S_i\cup F_i$ with $h_u=\sigma(S_i)$. Finally,
$F_i\setminus F\subseteq\{v_1,\dots,v_{i-1}\}\subseteq N(r)$, while $h_u=\sigma(S)+h_{v_i}$ has the entry $1+1=0$ in
the row $r$ (as $r\in U(S)$ and $v_i\in N(r)$), so $u\notin N(r)$ and $u\notin F_i\setminus F$. Hence the condition
$u\notin S_i\cup F_i$ is the same as $u\notin S\cup F\cup\{v_i\}$, for both values of $\varepsilon$. Two pairs
$(i,u)\neq(i',u')$ with $i\neq i'$ give different sets, since $\{v_i,u\}=\{v_{i'},u'\}$ would imply $u=v_{i'}\in N(r)$;
and pairs with $i=i'$ and $u\neq u'$ obviously give different sets. The case $|S|=W-1$ is the same argument for a
single level.
\end{proof}

\ttt{dist\_m4ri} finds the columns $u$ of~\eqref{eq:pair} without the branching over $A_{r'}(S_i,F_i)$: with random
values $z_1,\dots,z_m\in\{0,1\}^{64}$ and the hash $\zeta(x)=\bigoplus_{r\in\supp x}z_r$ of a vector $x\in\Ftwo^m$,
which satisfies $\zeta(x+y)=\zeta(x)\oplus\zeta(y)$, the nonzero columns of $H$ are stored by their hashes, and the
columns with $h_u=\sigma(S)+h_{v_i}$ are among those returned by a lookup of $\zeta(\sigma(S))\oplus\zeta(h_{v_i})$.
Every candidate is verified exactly, so that hash collisions cause no error, and equal vectors have equal hashes, so
that no column is missed. By Proposition~\ref{prop:hash}, the hash implementation reports the same sets as the search,
for both orders, and the condition $u\notin S\cup F\cup\{v_i\}$ is all that is needed for the forbidden set.

\begin{proposition}[work units]\label{prop:units}
Fix a start column and an integer $L\ge0$, and let $B\ge\max_r|N(r)|$, so that $q\le B$ at every node. For a vector of
digits $\delta\in[B]^L$, the \emph{unit} $\delta$ is the search from the start column in which every call at depth
$\ell-1<L$ (the depth of a call on $(S,F)$ is $|S|-1$) takes only the branch $i=\delta_\ell$, or none if
$\delta_\ell>q$, with $S_i$ and $F_i$ as in step~3(a); in particular, $v_1,\dots,v_{\delta_\ell-1}$ are forbidden in
it. Then every report event of the full search at a call of depth at least $L$ occurs in exactly one unit, and the
units have no other report events except at calls of depth less than $L$. With $\Wr=\{W\}$ and $L\le W-2$, all report
events are at calls of depth $W-2\ge L$, so that the units together report exactly the sets reported by the full
search, with the same multiplicities.
\end{proposition}

\begin{proof}
A call of the full search at depth $t$ is identified by its sequence of branch indices $(i_1,\dots,i_t)$. Since the
pruning, the choice rule, and the orderings are functions of the node, the unit $\delta$ makes exactly the calls of the
full search whose sequences are prefixes of $\delta$ (depth less than $L$) or begin with $\delta$ (depth at least $L$).
A report event at a call of depth $t\ge L$ therefore occurs in the unit $(i_1,\dots,i_L)$ and in no other unit.
\end{proof}

\ttt{dist\_m4ri} uses $B$ equal to the maximum row weight of $H$ and at most $L=\min(3,W-3)$ fixed levels, so that
the calls with $|S|=W-2$, where the hash table of Proposition~\ref{prop:hash} is used, are never at fixed levels.

\begin{remark}[threads and a decreasing weight limit]\label{rem:threads}
The search tree is determined by the nodes and the deterministic rules, so that the units can be processed by any
number of threads in any order. When a codeword is found, \ttt{dist\_m4ri} may lower the weight limit during a round
to the current upper bound plus \ttt{dW}. This only prunes more: for a codeword $C$ with $|C|$ at most the final limit,
every call on the branch of Theorem~\ref{thm:main} passes step~1 for every limit $W\ge|C|$, and its report event is
unique as before.
\end{remark}

\begin{thebibliography}{9}
\bibitem{DKP17}
I.~Dumer, A.~A.~Kovalev, and L.~P.~Pryadko, ``Distance verification for classical and quantum LDPC codes,''
\emph{IEEE Trans. Inf. Theory} \textbf{63}(7), 4675 (2017), doi:10.1109/TIT.2017.2690381, arXiv:1611.07164.
\bibitem{LOMKZ26}
J.~Y.~Lee, K.~Okada, N.~Maskara, K.~Kasai, and H.~Zhou, ``Design principles for ultra-high-rate quantum codes,''
arXiv:2609.30069 (2026), Appendix~B.
\bibitem{Lee26}
G.~Lee, ``Bicycle Race: small two-block codes on the Pareto frontier,'' arXiv:2610.13880 (2026), Appendix~C.4.
\end{thebibliography}

\end{document}
